\documentclass[11pt]{article}
\title{From Warp Metrics to Engineering Requirements: Conditional Source Constraints and an Operational Travel Benchmark}
\date{}
\author{SIGMA AI-assisted research synthesis}
\usepackage[T1]{fontenc}
\usepackage[utf8]{inputenc}
\usepackage[margin=1in]{geometry}
\usepackage{amsmath,amssymb,booktabs,longtable}
\usepackage{microtype}
\usepackage{titlesec}
\emergencystretch=2em
\raggedbottom
\linespread{1.07}
\usepackage{xcolor}
\definecolor{blue}{HTML}{24485E}
\titleformat{\section}{\Large\bfseries\sffamily\color{blue}}{\thesection}{0.8em}{}
\titleformat{\subsection}{\large\bfseries\sffamily}{\thesubsection}{0.7em}{}
\setlength{\parskip}{.6em}
\setlength{\parindent}{0pt}
\begin{document}
\maketitle
\begin{center}\small\textit{AI-generated research synthesis \textperiodcentered{} SIGMA}\end{center}

\begin{abstract}
Can a warp-like spacetime geometry become a propulsion system without requiring faster-than-light travel? This paper develops a bounded assessment of explicitly prescribed, flat-slice geometries in classical general relativity. The strongest analytic result is a conditional source obstruction: a smooth spherical shift bubble requires negative energy density for slice-normal observers at every nonzero speed, including subluminal speeds. A profile-independent lower bound quantifies the magnitude of its negative slice-energy integral, but does not determine fuel consumption or construction work. A separate subluminal benchmark distinguishes local motion, traveler aging, and synchronized laboratory arrival times.

Its center rider has zero local velocity relative to the normal observers, yet arrives later than both a selected light signal in the actual geometry and a photon in a matched Minkowski experiment. These results remain unverified analytic candidates, not demonstrated propulsion performance. They identify concrete source-compatibility and operational tests while leaving material realization, formation, stability, control, and navigation unresolved.
\end{abstract}
\section*{About this paper}
This paper was generated by SIGMA, an AI-assisted research system, from the project's recorded investigations, findings, and supporting evidence.
The analysis and conclusions reflect an automated synthesis of that research. Although SIGMA incorporates evidence tracking, critical review, and consistency checks, the resulting paper has not necessarily undergone independent verification or peer review.
Readers should evaluate its arguments, supporting evidence, and stated limitations before relying on its conclusions.
Research style: Distinguish simulated, measured and predicted performance, with operating constraints, failure modes and validation gaps.
\section*{Evidence coverage limitation}
The evidence index omitted 1152 material saved research field fragment(s) across 3 source(s). These fragments were not reviewed for this paper. Interpretations that could depend on omitted material remain provisional; the complete frozen source snapshot and omission ledger are retained with the generation record.
\textit{Claims discussed: }C12, C16
\section{Introduction and design objectives}
The engineering question is whether a prescribed geometry can be supported, created, controlled, and used under explicit operating constraints—not merely whether its line element admits a moving timelike center. The available analysis addresses two specified flat-slice constructions rather than the full range of warp-like proposals. It therefore supports candidate-specific failure tests and operational comparisons, not a universal verdict on spacetime-based propulsion.

Subluminal operation separates two questions often conflated in warp-drive discussions: whether a traveler moves locally faster than light, and whether a required source can physically exist. The assessment treats observer-measured energy density, laboratory travel time, traveler proper time, and global causal structure as distinct quantities. Each has its own assumptions and acceptance criteria.

\textit{Claims discussed: }C10, C12
\section{Methods and evidence}
The source analysis adopted classical Einstein gravity with zero cosmological constant and interpreted the Einstein tensor as the total required stress-energy. For a unit-lapse metric with Euclidean spatial slices and a smooth radial shift profile, the saved calculation obtained the slice-normal energy density from the Hamiltonian constraint. Angular integration and weighted Cauchy–Schwarz then produced a profile-independent bound; dimensional analysis and rounded arithmetic supplied an illustrative scale.

The transport analysis prescribed a compact-support profile, a finite-duration subluminal center motion, and two exterior laboratories synchronized during an initially Minkowski interval. It calculated a piecewise-timelike traveler's laboratory and proper times, bounded a selected axial null signal's arrival, and compared the result with a separately matched Minkowski experiment. A space-dependent change of time coordinate checked whether the operational timestamp differences survived coordinate relabeling.

The engineering synthesis separated necessary source constraints from unavailable capabilities and compared the geometric construction with an ideal photon rocket governed by energy and momentum conservation. No experiment, numerical integration, or simulation was performed. Primary-source correspondence and independent verification of the analytic derivations remain outstanding.

\textit{Claims discussed: }C11, C12, C14, C16
\section{Geometry, observers, and the scope of an engineering claim}
A prescribed warp-like geometry becomes an engineering proposition only when its required stress-energy is distinguished from a physical system capable of producing it. The assessment here assumes classical Einstein gravity with zero cosmological constant and interprets $T_{ab}=c^4G_{ab}/(8\pi G)$ as the total required source, rather than as one contribution to an otherwise unspecified source. This relation determines what stress-energy a chosen geometry would require; it does not supply matter fields, material properties, or an operating mechanism that realizes that requirement.

The source-constraint model uses $x^0=ct$ and the line element $ds^2=-(dx^0)^2+[dx-\beta f(r_s)dx^0]^2+dy^2+dz^2$, where $r_s^2=(x-\beta x^0)^2+y^2+z^2$. The constant, dimensionless parameter $\beta$ specifies the center’s coordinate speed divided by $c$, so its prescribed position is $x_s=\beta ct$. Neither this displacement law nor the metric is adopted as evidence that a source can generate the motion.

The radial profile is smooth and satisfies $0\leq f\leq1$, with $f=1$ for $r_s\leq R$ and $f=0$ for $r_s\geq R+\Delta$, where $R>0$ and $\Delta>0$. Thus $R$ defines the inner constant-profile radius, while $\Delta$ defines the thickness of the intervening wall. The geometry has unit lapse, Euclidean spatial metric, and shift $N^i=(-\beta f,0,0)$. These are restrictive mathematical premises of the source calculation, not adjustable engineering subsystems.

Energy density must be assigned to specified observers. For this slicing, the future-directed unit normal is $n=\partial_{x^0}+\beta f\partial_x$, with $g(n,n)=-1$, and the corresponding measured energy density is $\varepsilon=T_{ab}n^an^b$. This contraction is expressed in energy-density units and refers to the slice-normal observers, rather than to an unspecified laboratory frame. Their timelike normalization makes them an admissible family of physical observers for testing the required source.

The weak energy condition requires $T_{ab}u^au^b\geq0$ for every timelike observer direction $u^a$. Consequently, a negative value of $\varepsilon$ at even one event would suffice to show that the required total source violates that condition; a complete reconstruction of every stress-tensor component is unnecessary for that exclusion. The total-source interpretation is essential: if $T_{ab}=\sum_A T^{(A)}_{ab}$ and every constituent satisfies the condition, their contractions with any fixed timelike direction add to a nonnegative value. Conversely, a nonnegative measurement for the normal observers alone would not establish the condition for all observers, much less demonstrate a realizable source.

The evidential scope is limited to two explicitly prescribed flat-slice constructions: this constant-speed radial-wall family and a separate compact-support adaptation used for an operational travel benchmark. The calculations do not constitute an exhaustive review or establish results for arbitrary lapse functions, curved spatial slices, alternative gravitational theories, or other propulsion geometries. Correspondence to Miguel Alcubierre’s identified paper, “The warp drive: hyper-fast travel within general relativity,” remains unverified because its primary text was not inspected; the equations and conclusions therefore apply to the constructions stated here, not automatically to that paper’s original formulation or travel-time claims.

The available results are conditional analytic candidates, accompanied by rounded arithmetic rather than executed numerical experiments, simulations, physical measurements, or independent formal verification. Their engineering significance is nevertheless precise: an observer-resolved source test can reject a proposed source class before a propulsion design exists, whereas specifying a metric cannot establish such a design. The relevant first question is therefore whether a candidate total source can satisfy the geometry’s requirements, not whether its prescribed coordinate displacement resembles successful transport.

\textit{Claims discussed: }C1, C2, C14, C16
\section{Negative energy at subluminal speed and its quantitative scale}
For the specified flat-slice bubble, the Hamiltonian constraint yields a conditional source obstruction that persists at subluminal speed, together with a profile-independent lower bound on the magnitude of its negative slice-energy integral. The argument assumes classical Einstein gravity with zero cosmological constant and interprets the required stress-energy as the total source of the prescribed geometry. It is an analytic candidate, not an independently validated source calculation or a demonstration of material realizability.

Write the shift as $N^i=(-\beta f,0,0)$, with constant dimensionless speed $\beta$, and denote spatial derivatives by $f_i=\partial_i f$. Because the spatial metric is the time-independent Euclidean metric, the convention $K_{ij}=(\partial_iN_j+\partial_jN_i)/2$ gives the nonzero extrinsic-curvature components
$K_{xx}=-\beta f_x,\qquad K_{xy}=K_{yx}=-\frac{\beta}{2}f_y,\qquad K_{xz}=K_{zx}=-\frac{\beta}{2}f_z.$
Their trace and quadratic contraction are therefore $K=-\beta f_x$ and $K_{ij}K^{ij}=\beta^2[f_x^2+(f_y^2+f_z^2)/2]$. Reversing the overall sign convention for extrinsic curvature leaves the quadratic combination used below unchanged.

The longitudinal-gradient contribution cancels exactly:
$K^2-K_{ij}K^{ij}=\beta^2f_x^2-\beta^2\left[f_x^2+\frac{f_y^2+f_z^2}{2}\right]=-\frac{\beta^2}{2}(f_y^2+f_z^2).$
The intrinsic spatial scalar curvature is ${}^{(3)}R=0$, so the Hamiltonian constraint $16\pi G\varepsilon/c^4={}^{(3)}R+K^2-K_{ij}K^{ij}$ fixes the slice-normal observer's energy density as
$\varepsilon=-\frac{c^4\beta^2}{32\pi G}(f_y^2+f_z^2).$
Thus the density is nonpositive everywhere, with strict negativity wherever a transverse profile gradient is nonzero; no positive intrinsic-curvature term is available in this geometry to offset it.

For the radial profile, let $\theta$ be measured from the translation axis, so that \$f\_y\textasciicircum{}2+f\_z\textasciicircum{}2=f'(r\_s)\textasciicircum{}2\textbackslash{}sin\textasciicircum{}2\textbackslash{}theta\$. The required transition from $f(R)=1$ to $f(R+\Delta)=0$ forces a nonzero derivative somewhere in the wall, and smoothness makes that derivative nonzero on an open radial interval. Consequently, whenever $\beta\neq0$, an open wall region away from the motion axis has $\varepsilon<0$, including for every $0<|\beta|<1$. An axis-only calculation would miss this obstruction because the transverse contribution vanishes there; the negative measurement by the timelike slice-normal observer suffices to violate the weak energy condition, without establishing violations of every other energy condition.

The spatial integral quantifies this necessary negative density without choosing a particular wall shape. Define $E_{\mathrm{slice}}=\int\varepsilon\,d^3x$ on a constant-time slice, using the Euclidean proper volume supplied by the spatial metric. Angular integration gives $\int\sin^2\theta\,d\Omega=8\pi/3$, and hence
\$E\_\{\textbackslash{}mathrm\{slice\}\}=-\textbackslash{}frac\{c\textasciicircum{}4\textbackslash{}beta\textasciicircum{}2\}\{12G\}I,\textbackslash{}qquad I=\textbackslash{}int\_R\textasciicircum{}\{R+\textbackslash{}Delta\}r\textasciicircum{}2f'(r)\textasciicircum{}2\textbackslash{},dr.\$
Here $r$ is the radial integration variable, and only the transition wall contributes because the profile is constant in the interior and exterior.

The fixed endpoint difference supplies a lower bound on $I$ through weighted Cauchy–Schwarz. Applying the inequality to \$rf'(r)\$ and $r^{-1}$ yields
\$1=\textbackslash{}left|\textbackslash{}int\_R\textasciicircum{}\{R+\textbackslash{}Delta\}f'(r)\textbackslash{},dr\textbackslash{}right|\textasciicircum{}2\textbackslash{}leq\textbackslash{}left[\textbackslash{}int\_R\textasciicircum{}\{R+\textbackslash{}Delta\}r\textasciicircum{}2f'(r)\textasciicircum{}2\textbackslash{},dr\textbackslash{}right]\textbackslash{}left[\textbackslash{}int\_R\textasciicircum{}\{R+\textbackslash{}Delta\}r\textasciicircum{}\{-2\}\textbackslash{},dr\textbackslash{}right]=I\textbackslash{}frac\{\textbackslash{}Delta\}\{R(R+\textbackslash{}Delta)\}.\$
It follows that
$|E_{\mathrm{slice}}|\geq\frac{c^4\beta^2}{12G}\frac{R(R+\Delta)}{\Delta}.$
The estimate uses the endpoint values rather than an assumed detailed interpolation, so it does not require monotonicity.

Equality in the integral inequality would require \$f'(r)\$ to be proportional to $r^{-2}$ throughout the wall. A nonzero proportionality constant would leave nonzero endpoint slopes, incompatible with smooth joins to both constant-profile regions, so attainment of the displayed bound is not claimed. Nor has a thin-wall approximation been used: at fixed $R>0$ and fixed nonzero $\beta$, the bound diverges as $\Delta\to0$. At $\beta=0$, both the density and the bound vanish, whereas reducing a nonzero speed decreases their scale quadratically without removing the negative region.

The dimensional check is consistent at both local and integrated levels. The factor $c^4/G$ has units $\mathrm{J/m}$, while a spatial profile derivative has units $\mathrm{m}^{-1}$, giving $\varepsilon$ units of $\mathrm{J/m^3}$. Both $I$ and $R(R+\Delta)/\Delta$ have units of length, so the slice integral and its lower bound have energy units.

For an illustrative subluminal choice $R=100\,\mathrm{m}$, $\Delta=10\,\mathrm{m}$, and $\beta=0.01$, the geometric-speed factor is $\beta^2R(R+\Delta)/\Delta=0.11\,\mathrm{m}$. Using $c\approx2.998\times10^8\,\mathrm{m/s}$ and $G\approx6.674\times10^{-11}\,\mathrm{m^3\,kg^{-1}\,s^{-2}}$ gives $c^4/(12G)\approx1.01\times10^{43}\,\mathrm{J/m}$ and therefore \$|E\_\{\textbackslash{}mathrm\{slice\}\}|\textbackslash{}gtrsim1.11\textbackslash{}times10\textasciicircum{}\{42\}\textbackslash{},\textbackslash{}mathrm\{J\}\$. This is rounded arithmetic conditional on the analytic bound, not measured or simulated propulsion performance.

The resulting quantity is an observer- and slicing-defined integral of negative local energy, not a conserved total energy, fuel requirement, construction-work estimate, or electrical-input budget. No source-production law or apparatus efficiency connects it to those engineering quantities, and adding sources that all satisfy the weak energy condition cannot reproduce the negative total density because their contractions with a fixed timelike observer add.

The evidence therefore shifts the feasibility question from whether this bubble travels below light speed to whether its required source is physically admissible: slowing this exact geometry changes the magnitude of the obstruction, not its sign, while a different geometry requires a new calculation rather than inheritance of this bound.

\textit{Claims discussed: }C3, C4, C5
\section{Laboratory travel times, local motion, and causal interpretation}
An operational travel benchmark separates three quantities that a prescribed warp-like geometry does not make interchangeable: local velocity, traveler proper time, and arrival time recorded by exterior laboratories. The compact-support construction considered here permits a rider to remain locally at rest relative to slice-normal observers while moving between separated laboratory locations, but its prescribed trip takes longer than either of two specified light-signal comparisons. These are conditional analytic results, not measured or numerically validated propulsion performance; they apply to the explicit adaptation below rather than to uninspected travel-time claims in the original Alcubierre paper.

Take the globally defined metric $ds^2=-c^2dt^2+[dx-v(t)f(r)dt]^2+dy^2+dz^2$ on $\mathbb{R}^4$, where $r^2=[x-X(t)]^2+y^2+z^2$ and $v=dX/dt$. For $R>0$, prescribe $f(r)=\exp[-r^2/(R^2-r^2)]$ when $r<R$, and $f(r)=0$ when $r\geq R$. This smooth profile satisfies $0\leq f\leq1$ and $f(0)=1$. Here $R$ is the outer support radius, not the inner radius of the constant-profile region used in the preceding wall model; the present profile does not contain that extended constant interior.

For $D>0$ and $s=t/T$ within $0\leq s\leq1$, prescribe the center motion as $X(t)=D(35s^4-84s^5+70s^6-20s^7)$, with $X=0$ before the interval and $X=D$ afterward. Differentiation gives $v(t)=(140D/T)s^3(1-s)^3$ during the interval and zero outside it. Since $s(1-s)$ reaches its maximum at $s=1/2$, the maximum coordinate speed is $v_{\max}=35D/(16T)$; choosing $T=4D/c$ therefore gives $v_{\max}=35c/64<c$. The joined center function is $C^3$, making the velocity and metric at least $C^2$, sufficient for the classical curvature calculation.

Place stationary laboratories A and B on the axis at $x_A=-L$ and $x_B=D+L$, with $L>R$. Because $0\leq X(t)\leq D$, their distances from the moving center never fall below $L$, so both remain outside the support throughout the disturbance. Before $t=0$, the shift vanishes everywhere and the spacetime is Minkowski: establish Einstein synchronization then, set both clock readings equal to $t$, and measure the initial separation as $S=D+2L$. Along both subsequent laboratory histories, $d\tau_A=d\tau_B=dt$, so their readings retain this relation without requiring a new synchronization procedure through the moving geometry.

Define one complete idealized trip using a fixed exterior-leg speed $0<u<c$. The traveler departs A at $E=(-L/u,-L,0,0)$, reaches the center at $(t,x)=(0,0)$, follows $x=X(t)$ until $(t,x)=(T,D)$, and then proceeds to B at speed $u$. Both connecting legs lie in globally flat time intervals, and reception occurs at $P=(T+L/u,D+L,0,0)$. Consequently, B's reception reading minus A's departure reading is $\Delta_{\mathrm{lab}}=T+2L/u$; this is a synchronized cross-laboratory observable, not a single clock's elapsed proper time.

Local motion requires a different calculation. The orthonormal coframe of the slice-normal observers contains the temporal differential $c\,dt$ and the axial spatial differential $dx-vf\,dt$, so these observers measure an axial velocity $U=dx/dt-vf$. At the center, $dx/dt=v$ and $f=1$, giving $U=0$ and $ds^2=-c^2dt^2$: the center rider accumulates proper time $T$ while the center moves through coordinate displacement $D$. Null motion instead satisfies \$dx/dt-vf=\textbackslash{}pm c\$, so the distinction between displacement across the geometry and local speed does not alter the locally measured speed of light.

Each flat connecting leg contributes proper time $(L/u)\sqrt{1-u^2/c^2}$, yielding the complete wristwatch reading $\tau_{\mathrm{trip}}=T+(2L/u)\sqrt{1-u^2/c^2}$. Thus the center segment contributes equally to laboratory duration and traveler aging, whereas the connecting legs contribute differently to these observables. The dimensions are consistent: $f$ is dimensionless, $vf$ has units of speed, and $T$, $L/u$, and the resulting clock differences have units of time. Removing the shift while retaining the same moving center path would change its proper-time contribution to $\int_0^T\sqrt{1-v(t)^2/c^2}\,dt$, rather than leave it equal to $T$; the aging result depends on the specified geometry, not merely on the drawn trajectory.

For the first arrival comparison, emit a forward axial light signal from A at the same departure event $E$. Reflection symmetry permits propagation along the axis, where the outgoing null equation during the disturbance is $dx_\gamma/dt=c+vf$. Because $v\geq0$ and $f\geq0$, its forward coordinate speed is at least $c$, and it equals $c$ in the flat intervals; integration over the laboratory separation therefore gives $t_\gamma-t_E\leq S/c$. This bound also covers the case in which the photon reaches B before the disturbance begins.

If the photon is still in transit at $t=0$, its flat incoming propagation has already placed it ahead of the center, at $x_\gamma(0)=L(c/u-1)>0$. During the disturbance its lead $q=x_\gamma-X$ satisfies $dq/dt=c-v(1-f)\geq c-v_{\max}>0$, so the center cannot overtake it. Neither this argument nor the arrival bound requires the selected axial ray to be the globally earliest causal signal. Any earlier causal reception would only strengthen the comparison with the traveler.

The second comparison is a separate Minkowski experiment with the same initial separation $S$, stationary laboratory histories, initial clock synchronization, and departure reading. Its forward photon takes $S/c$, whereas the prescribed traveler satisfies $\Delta_{\mathrm{lab}}-S/c=3D/c+2L(1/u-1/c)>0$ for $T=4D/c$. Combining this strict inequality with the actual-geometry photon bound shows that the traveler arrives later than both signals. The reference is explicit but not universal: matching initial separation and clock readings does not uniquely prescribe every alternative experiment or every vehicle comparison.

The clock comparison survives an explicit change of time coordinates. Under \$t'=t+\textbackslash{}eta x/c\$ and \$x'=x\$, with $\eta=1/4$, a raw coordinate difference between departure at A and reception at B becomes \$\textbackslash{}Delta t'=\textbackslash{}Delta t+\textbackslash{}eta S/c\$. The initially synchronized stationary laboratory clocks instead read \$\textbackslash{}tau=t'-\textbackslash{}eta x\_\{\textbackslash{}mathrm\{lab\}\}/c\$, so their operational timestamp difference is \$\textbackslash{}Delta t'-\textbackslash{}eta S/c=\textbackslash{}Delta t\$. On the center path, \$dt'=(1+\textbackslash{}eta v/c)dt\$ and \$d\textbackslash{}tau=dt=dt'/(1+\textbackslash{}eta v/c)\$, preserving the traveler proper time as well as the laboratory comparison.

Causal interpretation likewise depends on the complete spacetime rather than a coordinate speed alone. For this globally defined unit-lapse metric, $g^{ab}(\partial_a t)(\partial_b t)=-1/c^2$, so the gradient of $t$ is timelike. With the stipulated time orientation, $t$ increases along future-directed causal curves and therefore excludes closed causal curves within this domain. This result does not establish causal actuation or safe control, and it does not transfer automatically to combined moving bubbles, altered identifications, or different extensions.

The remaining trajectory idealization matters directly to engineering interpretation: joining the constant-speed connecting legs to the center ride introduces instantaneous velocity changes at $t=0$ and $t=T$. Although the piecewise timelike curve has the proper time calculated above, those corners require unbounded acceleration; bounded-proper-acceleration transfers and their arrival and aging costs remain to be calculated. Nor does losing to a photon determine performance against a vehicle constrained by specified thrust, energy, aging, or infrastructure, because no such resource-matched vehicle has been defined.

The benchmark therefore removes zero slice-normal velocity as evidence of a travel-time advantage: for this construction, local rest coexists with a strictly later operational arrival, leaving practical transport performance to be established by a finite-acceleration, explicitly constrained mission comparison.

\textit{Claims discussed: }C6, C7, C8, C9, C10, C15
\section{Engineering requirements and a momentum-exchange alternative}
The engineering consequence of the source obstruction is a set of acceptance requirements, not a propulsion specification. A useful requirements matrix must distinguish framework requirements, such as matching the complete stress tensor and satisfying source dynamics; model-dependent deductions, such as the negative density required by the selected geometry; and capabilities for which no implementation has been supplied, including formation, control, and navigation. The missing capabilities are unresolved deliverables rather than observed subsystem failures, and the supplied results provide analytic constraints rather than measured or simulated engineering performance.

For stress-energy, the framework requirement is to reproduce all components of $T_{ab}=c^4G_{ab}/(8\pi G)$, including momentum density and spatial stresses, with a source that also satisfies its own equations of motion. The model-dependent density result already excludes an entirely weak-energy-condition-respecting total source for this exact geometry, even at subluminal speed. Adding constituents that each satisfy that condition cannot repair the mismatch, because their measured energy densities add for any fixed timelike observer. The missing deliverable is therefore a compatible complete source, not an inventory of positive stored energy or an apparatus whose nominal energy equals the magnitude of the negative slice integral.

The material-system requirement is distinct from the tensor calculation. A candidate must specify constitutive relations, attainable stresses, operating temperature, losses, spatial configuration, and lifetime, and show that these limits accommodate the required source throughout operation. No such material system is established by the supplied analysis. In particular, neither the sign nor the magnitude of the slice-energy integral supplies a source-production law, conversion efficiency, power budget, or construction work estimate; these quantities cannot be inferred by relabeling that integral as fuel.

Formation and shutdown require an evolution from stated, admissible initial data rather than a sequence of prescribed metric configurations. The acceptance deliverable is a source-evolution model that forms, accelerates, steers, and removes the disturbance while preserving the gravitational constraints. Choosing a center trajectory $X(t)$ specifies desired motion but does not identify the actuators or source changes that produce it. No formation, steering, or shutdown protocol has been supplied, so their feasibility cannot be inferred from the regularity of the prescribed trajectory.

Stability is a coupled source–geometry question, not a consequence of metric smoothness. Its acceptance test must evolve metric perturbations together with perturbations of the proposed dynamical source, satisfy the constraints, and compare any growth with the required mission duration. A smooth background supplies neither that evolution nor an allowable disturbance margin. The present evidence therefore identifies a necessary stability calculation without establishing either stable operation or an instability of a specified physical implementation.

Causal actuation adds an operational requirement beyond the causal properties of a fixed spacetime. A proposed actuator and source-evolution model must support a well-defined causal initial-value problem through startup, acceleration, steering, and shutdown; broader configurations would require their own assessment. The temporal-function result for the stipulated single geometry does not establish that its source can be commanded in this way. The missing deliverable is an explicit causal control protocol, not merely a chosen coordinate motion, and the analysis supplies no basis for replacing that protocol with either an assumed control capability or a coordinate-speed paradox.

Navigation and payload exposure require quantitative calculations along explicit trajectories in a specified external environment. The navigation deliverables include operational clock comparisons, photon propagation, sensing delays, and docking constraints; the exposure deliverables include payload and wall tidal tensors and collision and radiation loads. These calculations must be compared with stated acceleration, material, and payload tolerances before safety margins can be assigned. A timelike center trajectory alone establishes none of those margins, and no quantitative navigation or payload-safety envelope has been supplied.

The transport benchmark also leaves a specific trajectory defect that matters independently of source realizability. Its constant-speed exterior legs join the center ride through instantaneous velocity changes at boarding and exit: the resulting piecewise timelike curve has a well-defined proper time, but the corners require unbounded acceleration. A finite-acceleration mission assessment must replace them with specified bounded-proper-acceleration transfers and recalculate arrival times and traveler aging. This is a concrete refinement of the operational benchmark, not a solution to formation or control of the geometry.

A practical performance comparison must additionally specify the reference vehicle's thrust, energy, traveler-aging, and infrastructure constraints, and include preparation time and source infrastructure on the geometry side. Losing to a photon does not determine performance against a slower vehicle subject to those constraints. Conversely, no advantage over such a vehicle follows when its resource limits have not been defined. The existing clock comparison is therefore not a resource-matched mission comparison, and the slice-energy integral cannot supply the missing common resource accounting.

An ideal photon rocket provides a contrasting model in which the acceleration mechanism is explicitly momentum exchange rather than a prescribed spacetime shift. Assume an initially stationary rocket, perfect conversion into perfectly collimated backward photons, and a single burn ending at speed $bc$, where $0<b<1$ and $\gamma=(1-b^2)^{-1/2}$. If $M_i$ and $M_f$ are the initial and final rest masses and $E_\gamma$ is the emitted photon energy in the initial rest frame, conservation gives $M_ic^2=\gamma M_fc^2+E_\gamma$ and $E_\gamma/c=\gamma M_fbc$. Eliminating $E_\gamma$ yields $M_i/M_f=\gamma(1+b)=\sqrt{(1+b)/(1-b)}$, which is finite at every strictly subluminal target speed; at $b=0.1$, the ratio is approximately $1.106$.

This photon-rocket result is a conservation-law idealization, not demonstrated propulsion performance or a complete mission calculation. Positive-energy radiation and ordinary positive-density matter need not encounter the selected negative-density obstruction, but conversion efficiency, tolerable power density, collimation, beam handling, heat rejection, navigation, and stopping remain unspecified. Its single-burn mass ratio and the bubble's observer-defined slice-energy integral are not interchangeable resource measures.

The comparison changes the engineering interpretation of the warp-like candidate: subluminal transport does not itself demand the obstructed source, so the proposed geometry must establish both a compatible dynamical implementation and an operational benefit under stated constraints, rather than treating mathematical transport as that benefit.

\textit{Claims discussed: }C11, C12, C15
\section{Source compatibility as the next feasibility test}
Source compatibility is the next discriminating feasibility test: a prescribed geometry must be matched to a specified physical source, not merely assigned an effective stress tensor. The immediate priority is therefore to calculate the complete source required by an explicit smooth wall profile and determine whether a candidate classical or semiclassical theory can reproduce it from admissible initial data. This would turn the present necessary-source constraint into a test of a particular geometry–source combination, without presuming that a compatible source exists.

The first calculation should derive the full Einstein tensor and express $T_{ab}=c^{4}G_{ab}/(8\pi G)$ in the slice-normal orthonormal frame, retaining the stated classical Einstein equation and zero cosmological constant. It should include energy density, momentum density, energy fluxes, and spatial stresses, with the normal-normal component checked against the Hamiltonian-constraint result. Agreement with the density alone would leave the remaining source requirements undetermined; the fluxes and principal stresses may impose restrictions absent from that single contraction.

No complete orthonormal tensor has yet been supplied, and the existing evidence consists of conditional analytic derivations and rounded arithmetic, not executed simulations, physical measurements, or independent formal verification.

The distinction between an effective tensor and a material source is decisive at this stage. Conservation of the effective tensor through the Bianchi identity is a necessary consistency property, but it does not exhibit matter fields or constitutive relations capable of producing that tensor. A candidate source must satisfy its own equations of motion and initial-data constraints while reproducing all required stress components, including the negative normal-observer density. The useful outcome is either an explicit compatible solution or a demonstrated incompatibility for that specified source class; formal conservation alone supplies neither.

For a quantum source, the calculation must begin with explicit quantum fields, a state, boundary conditions, the background curvature, and sampling worldlines. These specifications are needed before calculating the full renormalized stress tensor and asking whether its magnitude, spatial distribution, duration, fluxes, and stresses match the geometric requirement. Exhibiting a negative expectation value somewhere would not establish that match. Nor does the classical density obstruction, by itself, validate or exclude quantum sourcing.

Quantum energy inequalities must then be applied only with their field, state-class, curvature, and sampling assumptions established from the relevant primary sources. No quantitative quantum bound has been established for the present construction, so a bound derived under different assumptions cannot simply be transferred to its wall. A candidate must also pass a self-consistent semiclassical backreaction calculation: matching a stress expectation value on an imposed background is not yet evidence that the coupled source and geometry sustain one another. The adequacy of the proposed state’s duration and fluctuations likewise remains part of the compatibility question rather than an available engineering capability.

Changing the geometry is a separate route that requires a separate source calculation. The supplied assessment covers two explicitly prescribed flat-slice constructions; it does not establish a feasibility theorem for arbitrary lapse functions, curved spatial slices, alternative shifts, supporting infrastructure, or different gravitational theories. For any revision, the discriminating calculation is a new Einstein-tensor and energy-condition analysis, not an unchanged application of the radial flat-slice bound. In particular, nonnegative energy density for one observer cannot establish the weak energy condition, which requires nonnegative density for every timelike observer.

A source model that survives these tests would give a physical basis for studying formation and coupled stability. The next step would be to evolve admissible source–geometry initial data and perturbations, rather than infer stability from the smoothness of a prescribed metric or infer control from a chosen center trajectory. Constraint-preserving formation, steering, and shutdown would then have to be assessed within that same dynamical model. Preliminary geometric hazard calculations can proceed before source compatibility is established, but they cannot substitute for a source evolution that actually produces the proposed motion.

The operational benchmark has a complementary, narrower refinement available without resolving sourcing. Its instantaneous boarding and exit velocity changes should be replaced by transfers with specified bounded proper acceleration, followed by recalculation of traveler proper time and receiver-clock arrival margins. This would remove an unbounded-acceleration idealization, but would not supply formation or actuation of the geometry. Any subsequent claim of advantage over a practical vehicle also requires explicit reference constraints on thrust, energy, traveler aging, and infrastructure; losing to a photon does not settle that resource-constrained comparison.

Numerical integration of the stated outgoing null equation at several geometric scale ratios, with convergence checks against the analytic arrival bound, could quantify the selected signal’s advance relative to the matched flat reference. That calculation remains unperformed, and its purpose would be to check implementation and sharpen the timing prediction, not demonstrate source realizability. Separately, the identified Alcubierre primary text must be inspected before claiming exact equation correspondence or attributing the present timing conclusions to its original construction; the current results apply to the explicitly stated geometries.

The evidence thus shifts the engineering question from choosing a coordinate speed to solving a constrained source problem. A successful full-tensor match, with admissible source dynamics and self-consistent backreaction where applicable, would justify advancing a particular candidate to coupled stability and control analysis; an incompatibility would reject that geometry–source pairing. Neither outcome would decide every warp-like proposal, but each would discriminate physical candidates more sharply than prescribing another trajectory alone.

\textit{Claims discussed: }C1, C12, C13, C14, C15, C16
\section{Design analysis and tradeoffs}
\subsection{C1 -- Limitation}
The available evidence supports a bounded assessment of two explicitly prescribed flat-slice geometries, not an exhaustive review or a universal feasibility theorem for warp-like propulsion.
The source constraint assumes one unit-lapse, Euclidean-slice radial shift family. The travel calculation uses a separately prescribed compact-support adaptation and trajectory. Neither calculation supplies results for arbitrary lapse functions, spatial curvature, alternative gravitational theories, or other proposed propulsion geometries.
\textit{Evidence: }artifact:745ebe9c-6fd0-4f37-8dc9-007be29c954c:eedd3fa0-719b-4b64-9277-bb3a21cf5284, artifact:745ebe9c-6fd0-4f37-8dc9-007be29c954c:027a5b69-a96f-43f7-91ce-60af8ce26b57
\subsection{C2 -- Established}
The source-bound model is defined by $x^0=ct$, $ds^2=-(dx^0)^2+[dx-\beta f(r_s)dx^0]^2+dy^2+dz^2$, and $r_s^2=(x-\beta x^0)^2+y^2+z^2$, with constant $\beta$, smooth $0\leq f\leq1$, $f=1$ for $r_s\leq R$, and $f=0$ for $r_s\geq R+\Delta$, where $R>0$ and $\Delta>0$.
This is an explicit mathematical specification, not a claim of physical realizability. Its lapse is one, its spatial metric is Euclidean, and its shift is $N^i=(-\beta f,0,0)$. The future unit normal is $n=\partial_{x^0}+\beta f\partial_x$, with $g(n,n)=-1$; the observer-measured energy density is $\varepsilon=T_{ab}n^an^b$.
\textit{Evidence: }artifact:745ebe9c-6fd0-4f37-8dc9-007be29c954c:eedd3fa0-719b-4b64-9277-bb3a21cf5284
\subsection{C3 -- Conjecture}
Conditional on Einstein gravity with zero cosmological constant and the specified total source interpretation, the source-bound model requires $\varepsilon=-c^4\beta^2(f_y^2+f_z^2)/(32\pi G)$ and violates the weak energy condition somewhere in the wall whenever $\beta\neq0$, including $0<|\beta|<1$.
This status denotes a supported but independently unverified analytic derivation. With $K_{ij}=(\partial_iN_j+\partial_jN_i)/2$, the nonzero entries are $K_{xx}=-\beta f_x$, $K_{xy}=K_{yx}=-\beta f_y/2$, and $K_{xz}=K_{zx}=-\beta f_z/2$. Consequently $K^2-K_{ij}K^{ij}=-\beta^2(f_y^2+f_z^2)/2$; the spatial scalar curvature vanishes, and the Hamiltonian constraint gives the stated density. For a radial profile, \$f\_y\textasciicircum{}2+f\_z\textasciicircum{}2=f'(r\_s)\textasciicircum{}2\textbackslash{}sin\textasciicircum{}2\textbackslash{}theta\$; the transition from one to zero requires a nonzero derivative on an open region, so negative density occurs away from the translation axis. The weak energy condition requires $T_{ab}u^au^b\geq0$ for every timelike observer, making this one negative observer measurement sufficient to exclude an everywhere-condition-respecting total source.
\textit{Evidence: }artifact:745ebe9c-6fd0-4f37-8dc9-007be29c954c:eedd3fa0-719b-4b64-9277-bb3a21cf5284, session:745ebe9c-6fd0-4f37-8dc9-007be29c954c
\subsection{C4 -- Conjecture}
For the source-bound model, $E_{\mathrm{slice}}=\int\varepsilon\,d^3x$ satisfies $|E_{\mathrm{slice}}|\geq[c^4\beta^2/(12G)]R(R+\Delta)/\Delta$ without a monotonicity assumption or thin-wall approximation.
The saved proof candidate uses the Euclidean proper volume and $\int\sin^2\theta\,d\Omega=8\pi/3$ to obtain $E_{\mathrm{slice}}=-[c^4\beta^2/(12G)]I$, where \$I=\textbackslash{}int\_R\textasciicircum{}\{R+\textbackslash{}Delta\}r\textasciicircum{}2f'(r)\textasciicircum{}2\textbackslash{},dr\$. The endpoint difference gives \$1=|\textbackslash{}int\_R\textasciicircum{}\{R+\textbackslash{}Delta\}f'\textbackslash{},dr|\textasciicircum{}2\textbackslash{}leq I\textbackslash{}int\_R\textasciicircum{}\{R+\textbackslash{}Delta\}r\textasciicircum{}\{-2\}\textbackslash{},dr=I\textbackslash{}Delta/[R(R+\textbackslash{}Delta)]\$. Equality would require \$f'\$ proportional to $r^{-2}$ throughout the wall, which cannot join smoothly to both constant regions, so attainment is not claimed. At $\beta=0$ both density and bound vanish; at fixed nonzero speed and fixed $R$, the bound diverges as $\Delta\to0$.
\textit{Evidence: }artifact:745ebe9c-6fd0-4f37-8dc9-007be29c954c:eedd3fa0-719b-4b64-9277-bb3a21cf5284, session:745ebe9c-6fd0-4f37-8dc9-007be29c954c
\subsection{C5 -- Established}
The slice-energy formulas are dimensionally consistent, and the illustrative substitution $R=100\,\mathrm{m}$, $\Delta=10\,\mathrm{m}$, and $\beta=0.01$ gives approximately \$|E\_\{\textbackslash{}mathrm\{slice\}\}|\textbackslash{}gtrsim1.11\textbackslash{}times10\textasciicircum{}\{42\}\textbackslash{},\textbackslash{}mathrm\{J\}\$; this is not a fuel requirement, construction-work estimate, or conserved total energy.
The factor $c^4/G$ has units $\mathrm{J/m}$, a spatial profile derivative has units $\mathrm{m}^{-1}$, and $R(R+\Delta)/\Delta$ has units of length. The substitution therefore has energy units and is rounded arithmetic conditional on the analytic bound, not measured or simulated performance. The integral is tied to the specified observers and slicing; no source-production law, apparatus efficiency, or creation protocol connects it to usable stored energy or input power.
\textit{Evidence: }artifact:745ebe9c-6fd0-4f37-8dc9-007be29c954c:eedd3fa0-719b-4b64-9277-bb3a21cf5284, session:745ebe9c-6fd0-4f37-8dc9-007be29c954c
\subsection{C6 -- Established}
The operational benchmark prescribes a compact-support radial profile and a finite-duration center trajectory whose maximum coordinate speed is $35c/64$, with two laboratories that remain outside the disturbance and retain an initially established clock synchronization.
The premises are $ds^2=-c^2dt^2+[dx-v(t)f(r)dt]^2+dy^2+dz^2$, $r^2=[x-X(t)]^2+y^2+z^2$, and $f=\exp[-r^2/(R^2-r^2)]$ for $r<R$, with $f=0$ otherwise. For $s=t/T$ in $[0,1]$, $X=D(35s^4-84s^5+70s^6-20s^7)$; it equals zero before the interval and $D$ afterward. Thus $v=(140D/T)s^3(1-s)^3$ and, for $T=4D/c$, $v_{\max}=35c/64<c$. Laboratories at $x=-L$ and $x=D+L$, with $L>R$, remain exterior; their clocks satisfy $d\tau=dt$, and their initial measured separation is $S=D+2L$.
\textit{Evidence: }artifact:745ebe9c-6fd0-4f37-8dc9-007be29c954c:027a5b69-a96f-43f7-91ce-60af8ce26b57
\subsection{C7 -- Conjecture}
In the operational benchmark, a center rider has zero local velocity relative to the slice-normal observers, while the complete idealized trip has synchronized laboratory duration $\Delta_{\mathrm{lab}}=T+2L/u$ and traveler proper time $\tau_{\mathrm{trip}}=T+(2L/u)\sqrt{1-u^2/c^2}$ for $0<u<c$.
The saved analytic candidate joins a flat incoming leg from $(-L/u,-L,0,0)$ to the center at $t=0$, a center segment $x=X(t)$, and a flat outgoing leg ending at $(T+L/u,D+L,0,0)$. Along the axis the normal observers measure $U=dx/dt-vf$, which vanishes at the center because $f=1$ and $dx/dt=v$. Substitution gives $ds^2=-c^2dt^2$ on that segment; the two exterior legs contribute the usual special-relativistic proper times. This distinguishes displacement between separated events from a locally measured transport speed.
\textit{Evidence: }artifact:745ebe9c-6fd0-4f37-8dc9-007be29c954c:027a5b69-a96f-43f7-91ce-60af8ce26b57
\subsection{C8 -- Conjecture}
The benchmark traveler arrives later than a selected forward light signal in the actual geometry and later than a photon in the specified matched Minkowski reference experiment.
During the disturbance the selected axial null signal obeys $dx_\gamma/dt=c+vf\geq c$, and in the flat intervals its forward speed is $c$, giving $t_\gamma-t_E\leq S/c$. If still in transit, its lead over the center satisfies $d(x_\gamma-X)/dt=c-v(1-f)\geq c-v_{\max}>0$. The reference is a separate Minkowski experiment with the same initial separation, clock synchronization, stationary laboratory histories, and departure reading, so its photon takes $S/c$. For $T=4D/c$, the traveler has $\Delta_{\mathrm{lab}}-S/c=3D/c+2L(1/u-1/c)>0$. The selected ray need not be globally earliest: any earlier causal signal only strengthens the comparison.
\textit{Evidence: }artifact:745ebe9c-6fd0-4f37-8dc9-007be29c954c:027a5b69-a96f-43f7-91ce-60af8ce26b57
\subsection{C9 -- Conjecture}
The saved coordinate-change check preserves the operational laboratory timestamp differences and traveler proper time under \$t'=t+\textbackslash{}eta x/c\$, \$x'=x\$, with $\eta=1/4$.
Raw cross-laboratory coordinate differences acquire $\eta S/c$, but a stationary laboratory's synchronized clock reads \$\textbackslash{}tau=t'-\textbackslash{}eta x\_\{\textbackslash{}mathrm\{lab\}\}/c\$. Subtracting these fixed clock offsets recovers the original operational difference. On the center curve, \$dt'=(1+\textbackslash{}eta v/c)dt\$ and \$d\textbackslash{}tau=dt=dt'/(1+\textbackslash{}eta v/c)\$. This checks the stated comparison under one explicit coordinate transformation; it does not uniquely define every possible reference experiment.
\textit{Evidence: }artifact:745ebe9c-6fd0-4f37-8dc9-007be29c954c:027a5b69-a96f-43f7-91ce-60af8ce26b57
\subsection{C10 -- Conjecture}
For the explicitly globally defined unit-lapse metric, the temporal-function argument excludes closed causal curves within the stated domain, without establishing safe control or causal consistency of different multi-bubble constructions.
The saved derivation gives $g^{ab}(\partial_at)(\partial_bt)=-1/c^2$, so the gradient of the globally defined time coordinate is timelike. With the stated time orientation, this supplies a temporal function and rules out a closed causal curve within that single spacetime. The result does not transfer automatically to altered identifications or combined moving configurations. Coordinate displacement exceeding a reference light-travel distance would not, by itself, establish a local speed violation or a causal paradox.
\textit{Evidence: }artifact:745ebe9c-6fd0-4f37-8dc9-007be29c954c:eedd3fa0-719b-4b64-9277-bb3a21cf5284
\subsection{C11 -- Conjecture}
An ideal single-burn photon rocket provides a contrasting momentum-exchange model with $M_i/M_f=\sqrt{(1+b)/(1-b)}$ for $0<b<1$, avoiding the selected negative-density obstruction without demonstrating practical propulsion capability.
The saved analytic comparison assumes an initially stationary rocket and perfectly collimated backward photon emission. Energy and momentum conservation give $M_ic^2=\gamma M_fc^2+E_\gamma$ and $E_\gamma/c=\gamma M_fbc$, hence $M_i/M_f=\gamma(1+b)$. At $b=0.1$, the ratio is approximately $1.106$. This is a one-burn idealization, not a complete mission calculation; conversion efficiency, power density, heat rejection, beam handling, and stopping remain unspecified. Its mass ratio and the bubble's slice-energy integral are not interchangeable resource measures.
\textit{Evidence: }session:745ebe9c-6fd0-4f37-8dc9-007be29c954c
\subsection{C12 -- Unresolved}
No supplied result establishes a compatible material source, formation protocol, coupled stability, causal actuator system, steering or shutdown capability, or quantitative navigation and payload-safety margins.
The requirements analysis identifies these as missing deliverables rather than observed failures. A prescribed center trajectory is not a control law, and smoothness of the metric is not a stability calculation. Required acceptance tests include a full source tensor, a source model satisfying its own dynamics, admissible initial data, perturbation evolution over a mission duration, constitutive and thermal limits, and explicit sensing, tidal, radiation, collision, and docking calculations.
\textit{Evidence: }artifact:745ebe9c-6fd0-4f37-8dc9-007be29c954c:eedd3fa0-719b-4b64-9277-bb3a21cf5284, artifact:745ebe9c-6fd0-4f37-8dc9-007be29c954c:027a5b69-a96f-43f7-91ce-60af8ce26b57
\subsection{C13 -- Unresolved}
Quantum sourcing and revised geometries remain unresolved alternatives; neither is validated or excluded by the negative-density calculation for the selected classical metric family.
A quantum proposal must specify fields, state, boundaries, curvature, and sampling worldlines before calculating the renormalized stress tensor and applying relevant quantum energy inequalities. Merely producing some negative energy would not establish agreement with every required stress component or semiclassical backreaction. A changed geometry requires a new Einstein-tensor and energy-condition analysis; nonnegative density for one observer does not establish the weak energy condition for all observers.
\textit{Evidence: }session:745ebe9c-6fd0-4f37-8dc9-007be29c954c, artifact:745ebe9c-6fd0-4f37-8dc9-007be29c954c:eedd3fa0-719b-4b64-9277-bb3a21cf5284
\subsection{C14 -- Limitation}
Correspondence to the identified Alcubierre primary source remains unverified; the present calculations apply to the explicitly stated constructions rather than verified equations or travel-time claims from that paper.
The evidence supplies a bibliographic lead but no inspected primary text. Exact equation matches, conventions, original boundary behavior, and original quantitative performance claims therefore remain unchecked. The self-contained calculations can be discussed without treating that attribution as verified.
\textit{Evidence: }artifact:745ebe9c-6fd0-4f37-8dc9-007be29c954c:eedd3fa0-719b-4b64-9277-bb3a21cf5284, artifact:745ebe9c-6fd0-4f37-8dc9-007be29c954c:027a5b69-a96f-43f7-91ce-60af8ce26b57
\subsection{C15 -- Limitation}
The transport benchmark is not a finite-acceleration mission trajectory or a resource-matched comparison against a practical vehicle.
Joining constant-speed exterior legs to the center segment introduces instantaneous velocity changes at boarding and exit. The curve has a well-defined piecewise proper time, but those corners require unbounded acceleration and cannot represent physical transfer maneuvers. Losing to a photon does not settle performance against a vehicle constrained by thrust, energy, aging, or infrastructure. Those constraints have not been specified.
\textit{Evidence: }artifact:745ebe9c-6fd0-4f37-8dc9-007be29c954c:027a5b69-a96f-43f7-91ce-60af8ce26b57
\subsection{C16 -- Limitation}
The evidence contains analytic candidates and rounded arithmetic, but no executed numerical experiment, physical measurement, simulation result, or independent formal verification.
The saved assessments describe the derivations as conditional and unverified. Agreement with a derivation is not independent proof certification, and the illustrative energy scales are substitutions into formulas rather than apparatus observations. Proposed null-equation integration and convergence studies remain unperformed.
\textit{Evidence: }artifact:745ebe9c-6fd0-4f37-8dc9-007be29c954c:eedd3fa0-719b-4b64-9277-bb3a21cf5284, artifact:745ebe9c-6fd0-4f37-8dc9-007be29c954c:027a5b69-a96f-43f7-91ce-60af8ce26b57, session:745ebe9c-6fd0-4f37-8dc9-007be29c954c
\section{Limitations}
\begin{itemize}
\item The evidence is a bounded selection and does not establish exhaustive coverage of proposed warp geometries or the primary literature. Correspondence between the specified adaptations and the identified original source remains unchecked. [C1, C14]
\item The source analysis derives a necessary energy-density component and a slice integral, not the complete stress tensor, a source-production mechanism, or an energy-input budget. Its bound must not be transferred unchanged to different metric families or gravitational assumptions. [C3, C4, C5, C12, C13]
\item The travel calculation permits velocity corners and lacks a resource-matched reference vehicle. It establishes a conditional photon comparison, not a practical mission or a complete comparison of propulsion technologies. [C8, C15]
\item No measured, simulated, or experimentally demonstrated performance is available. The mathematical arguments remain independently unverified, while source realization and dynamical engineering requirements remain unresolved. [C12, C16]
\end{itemize}
\section{Open obligations}
\begin{itemize}
\item Inspect the identified primary source and compare defining equations, conventions, profile assumptions, and travel-time claims with the two stated constructions. [C14]
\item Independently reproduce the Hamiltonian-constraint calculation and derive the full Einstein tensor and orthonormal stress components for an explicit smooth wall profile. [C3, C4, C12, C16]
\item Specify a candidate material or field source and check its equations, admissible initial data, and backreaction against every required stress component. A conserved effective tensor alone does not satisfy this obligation. [C12, C13]
\item For any quantum proposal, specify fields, state, boundary conditions, curvature, and sampling worldlines; calculate the renormalized stress tensor and establish the applicability of relevant primary-source quantum bounds. [C13]
\item For any compatible source, analyze constraint-preserving formation, steering, shutdown, causal actuation, and coupled source–geometry perturbations over a stated mission duration. [C12]
\item Replace instantaneous boarding and exit velocity changes with bounded-proper-acceleration transfers and recompute arrival margins and traveler proper time. [C7, C15]
\item Define thrust, energy, aging, infrastructure, and stopping constraints for a reference vehicle before assessing a subluminal engineering advantage. [C11, C15]
\item Quantify tidal loads, material interfaces, radiation and collision exposure, sensing delays, navigation, and docking tolerances in an explicit environment. [C12]
\item Retain the proposed null-equation integration with $u=0.9c$, $L=2R$, several $D/R$ values, and convergence checks as an unperformed implementation study. Its purpose would be timestamp quantification, not demonstration of source realizability. [C8, C16]
\item Recompute the source constraints and energy-condition tests for any revised geometry rather than transferring the radial flat-slice bound. [C1, C13]
\end{itemize}
\section{Frozen obligation reconciliation}
Unlisted obligations are not presumed resolved. The archived research snapshot retains complete source references and the full reconciliation.
\begin{itemize}
\item Inspect the original Alcubierre paper and record exact equation and convention matches; the bibliographic lead is not a checked source extract. \textit{ (Not listed by author; not presumed resolved; 1 saved source reference(s))}
\item For a proposed quantum rescue, specify the quantum fields, state, boundary conditions, and sampling worldlines; calculate the full renormalized stress tensor and test both semiclassical Einstein consistency and applicable quantum energy inequalities. Merely exhibiting some negative energy is insufficient. \textit{ (Not listed by author; not presumed resolved; 1 saved source reference(s))}
\item The next discriminating calculation for any revised geometry is to derive its full Einstein tensor and evaluate its energy conditions, rather than transferring this flat-slice radial bound outside its assumptions. A nonnegative density for one observer does not prove the weak energy condition. \textit{ (Not listed by author; not presumed resolved; 1 saved source reference(s))}
\item No matter-field evolution, formation process, stability analysis, perturbative backreaction calculation, or causal control protocol has been supplied. A metric and its required source tensor do not establish these. \textit{ (Not listed by author; not presumed resolved; 1 saved source reference(s))}
\item A navigation and travel-time assessment still needs physical departure/arrival worldlines, clock comparisons, tidal-acceleration limits, and a specified external environment. Coordinate bubble speed alone is not that assessment. \textit{ (Not listed by author; not presumed resolved; 1 saved source reference(s))}
\item The ideal photon-rocket comparison establishes conservation-law consistency only. Real conversion efficiency, tolerable power density, collimation, thermal limits, and stopping requirements remain unresolved engineering constraints. \textit{ (Not listed by author; not presumed resolved; 1 saved source reference(s))}
\item Inspect the original Alcubierre paper to verify the bibliographic locator, conventions, and correspondence between its construction and the explicitly defined metric here. \textit{ (Not listed by author; not presumed resolved; 3 saved source reference(s))}
\item Independently check the full Einstein tensor and orthonormal stress components for an explicit smooth wall profile; the present calculation uses only the Hamiltonian constraint. \textit{ (Not listed by author; not presumed resolved; 3 saved source reference(s))}
\item Provide a dynamical source model, or prove its incompatibility, rather than treating a conserved effective stress tensor as an available material. \textit{ (Not listed by author; not presumed resolved; 3 saved source reference(s))}
\item For any quantum source proposal, identify primary-source bounds with their exact field, state, curvature, and sampling assumptions before transferring them to this geometry. \textit{ (Not listed by author; not presumed resolved; 3 saved source reference(s))}
\item Analyze constraint-preserving formation, shutdown, steering, and coupled source–geometry stability for any source model that survives compatibility testing. \textit{ (Not listed by author; not presumed resolved; 3 saved source reference(s))}
\item Develop operational travel-time, infrastructure, sensing, tidal-force, radiation, and collision benchmarks before making engineering-performance claims. \textit{ (Not listed by author; not presumed resolved; 3 saved source reference(s))}
\item Recalculate rather than transfer the negative-energy bound if the spatial slicing, metric family, gravitational theory, or supporting infrastructure is changed. \textit{ (Not listed by author; not presumed resolved; 3 saved source reference(s))}
\item Inspect the defining equations and travel-time statements in the identified Alcubierre primary source. Establish exactly which results transfer to its original profile and which depend on the compact-support adaptation used here. \textit{ (Not listed by author; not presumed resolved; 3 saved source reference(s))}
\item Independently check the derivation, especially the stress-energy convention and transformed clock bookkeeping; a symbolic calculation would be useful but would not itself certify a general theorem. \textit{ (Not listed by author; not presumed resolved; 2 saved source reference(s))}
\item Replace the idealized velocity corners with specified bounded proper accelerations and calculate the resulting proper times and arrival margins. \textit{ (Not listed by author; not presumed resolved; 3 saved source reference(s))}
\item If testing a claimed subluminal engineering advantage, specify the reference vehicle's allowed thrust, energy, aging, and infrastructure constraints before comparing it with the prescribed geometry. \textit{ (Not listed by author; not presumed resolved; 3 saved source reference(s))}
\item Calculate the complete required stress tensor and search for a specified matter or quantum-field model satisfying it, including conservation, backreaction, and admissible initial data. The negative normal-observer energy density cannot be omitted from that test. \textit{ (Not listed by author; not presumed resolved; 3 saved source reference(s))}
\end{itemize}
52 further distinct obligations are retained in the archived reconciliation.
\section{Conclusions}
Removing faster-than-light operation does not remove the negative-density requirement in the specified flat-slice spherical shift family, according to the supported but unverified analytic derivation. The result rejects an everywhere-weak-energy-condition-respecting total source for that candidate, not every warp-like geometry. [C1, C3]

The profile-independent slice-energy bound supplies a necessary source scale, not a fuel estimate. Dimensional consistency and a large illustrative magnitude do not establish a production mechanism or a feasible apparatus. [C4, C5]

Zero local velocity at the bubble center is compatible with displacement between exterior laboratories, but it does not establish a travel-time advantage. In the stated subluminal benchmark, the traveler loses to both the selected actual-geometry light signal and the matched reference photon, with the operational comparison preserved under the tested coordinate change. [C7, C8, C9]

The evidence neither demonstrates an engineering warp drive nor establishes a universal impossibility theorem. Full source compatibility is the next decision gate; stability, causal control, materials, navigation, safety, and resource-matched performance require separate quantitative demonstrations. [C1, C12, C13, C15, C16]

\textit{Claims discussed: }C1, C3, C4, C5, C7, C8, C9, C12, C13, C15, C16
\appendix
\section{Selected source-linked research excerpts}
The following are bounded, selected excerpts of original research records. They are not exhaustive; consult the full frozen snapshot archived with this paper. They do not constitute independent verification.
\subsection{Research excerpt 1}
\textit{Source and field: }artifact:745ebe9c-6fd0-4f37-8dc9-007be29c954c:eedd3fa0-719b-4b64-9277-bb3a21cf5284: skeptic\_output/reasoning\_assessments\par
[\{"step\_number": 1, "disposition": "accepted", "rationale": "Accepted as the explicit premises of the conditional claim, not as an established physical realization. Substitution gives g(n,n)=-1, and the lapse, shift, and Euclidean spatial metric follow from the line element. The primary-source citation is expressly identified as uninspected rather than used as proof.", "objection": "", "basis\_artifact\_ids": []\}, \{"step\_number": 2, "disposition": "accepted", "rationale": "The spatial metric is time independent even though the translating profile is not. Its shift derivatives give the listed K components; contraction cancels the longitudinal-gradient square and leaves the stated negative transverse contribution.

The Hamiltonian coefficient then gives ε=-c\textasciicircum{}4$\beta$²($f_y^{2}$+$f_z^{2}$)/(32$\pi$G). Smooth interpolation guarantees a nonzero derivative on an open radial interval, hence negative density off the motion axis for $\beta$$\neq$0. A negative contraction with a timelike normal suffices to violate the weak energy condition.", "objection": "", "basis\_artifact\_ids": []\}, \{"step\_number": 3, "disposition": "accepted", "rationale": "The angular factor is 8$\pi$/3, giving the coefficient 1/12 in the slice integral. Applying Cauchy–Schwarz to rf′ and 1/r yields I$\geq$R(R+$\Delta$)/$\Delta$. This uses only the endpoint values, not monotonicity. The equality profile would not join smoothly to both constant regions, but that does not weaken the non-strict lower bound.

The zero-speed and thin-wall limits follow directly.", "objection": "", "basis\_artifact\_ids": []\}, \{"step\_number": 4, "disposition": "accepted", "rationale": "c\textasciicircum{}4/G…

\subsection{Research excerpt 2}
\textit{Source and field: }session:745ebe9c-6fd0-4f37-8dc9-007be29c954c: steps/4/output/reasoning\_assessments\par
[\{"step\_number": 1, "disposition": "accepted", "rationale": "Accepted as the explicit premises of the conditional claim, not as an established physical realization. Substitution gives g(n,n)=-1, and the lapse, shift, and Euclidean spatial metric follow from the line element. The primary-source citation is expressly identified as uninspected rather than used as proof.", "objection": "", "basis\_artifact\_ids": []\}, \{"step\_number": 2, "disposition": "accepted", "rationale": "The spatial metric is time independent even though the translating profile is not. Its shift derivatives give the listed K components; contraction cancels the longitudinal-gradient square and leaves the stated negative transverse contribution.

The Hamiltonian coefficient then gives ε=-c\textasciicircum{}4$\beta$²($f_y^{2}$+$f_z^{2}$)/(32$\pi$G). Smooth interpolation guarantees a nonzero derivative on an open radial interval, hence negative density off the motion axis for $\beta$$\neq$0. A negative contraction with a timelike normal suffices to violate the weak energy condition.", "objection": "", "basis\_artifact\_ids": []\}, \{"step\_number": 3, "disposition": "accepted", "rationale": "The angular factor is 8$\pi$/3, giving the coefficient 1/12 in the slice integral. Applying Cauchy–Schwarz to rf′ and 1/r yields I$\geq$R(R+$\Delta$)/$\Delta$. This uses only the endpoint values, not monotonicity. The equality profile would not join smoothly to both constant regions, but that does not weaken the non-strict lower bound.

The zero-speed and thin-wall limits follow directly.", "objection": "", "basis\_artifact\_ids": []\}, \{"step\_number": 4, "disposition": "accepted", "rationale": "c\textasciicircum{}4/G…

\subsection{Research excerpt 3}
\textit{Source and field: }artifact:745ebe9c-6fd0-4f37-8dc9-007be29c954c:027a5b69-a96f-43f7-91ce-60af8ce26b57: skeptic\_output/remaining\_obligations\par
["Inspect the Alcubierre primary source and distinguish its equations and claims from the compact-support adaptation.", "Replace the velocity corners with bounded-proper-acceleration transfers and calculate their arrival and aging costs.", "Define resource and infrastructure constraints before claiming an advantage over a reference vehicle.", "Determine the full required stress tensor and test a specified classical or quantum source model, including admissible initial data and backreaction.", "Analyze creation, causal control, stability, tidal loads, materials, navigation, and shutdown for any proposed source.", "Retain the optional null-ODE convergence experiment as unperformed; its purpose is quantitative implementation checking, not proof of realizability.", "Inspect the defining equations and travel-time statements in the identified Alcubierre primary source.

Establish exactly which results transfer to its original profile and which depend on the compact-support adaptation used here.", "Replace the idealized velocity corners with specified bounded proper accelerations and calculate the resulting proper times and arrival margins.", "If testing a claimed subluminal engineering advantage, specify the reference vehicle's allowed thrust, energy, aging, and infrastructure constraints before comparing it with the prescribed geometry.", "Calculate the complete required stress tensor and search for a specified matter or quantum-field model satisfying it, including conservation, backreaction, and admissible initial data. The negative normal-observer energy density cannot be…

\subsection{Research excerpt 4}
\textit{Source and field: }artifact:745ebe9c-6fd0-4f37-8dc9-007be29c954c:eedd3fa0-719b-4b64-9277-bb3a21cf5284: skeptic\_output/proof\_obligations\par
["Local density sign: supported by cycle 1, investigator step 27d543af-13af-4f83-a2bc-f75edc5bad40, reasoning steps 1–2, using the explicit metric and Hamiltonian constraint.", "Profile-independent slice-energy bound: supported by cycle 1, investigator step 27d543af-13af-4f83-a2bc-f75edc5bad40, reasoning step 3, using angular integration and weighted Cauchy–Schwarz.", "Dimensional and illustrative arithmetic consistency: supported by cycle 1, investigator step 27d543af-13af-4f83-a2bc-f75edc5bad40, reasoning step 4; this is not executed numerical evidence.", "Local-motion and single-metric causality distinction: supported by cycle 1, investigator step 27d543af-13af-4f83-a2bc-f75edc5bad40, reasoning step 5, using the central line element and temporal function.", "Exclusion of entirely weak-energy-condition-respecting total sources: supported by cycle 1, investigator step 27d543af-13af-4f83-a2bc-f75edc5bad40, reasoning steps 2 and 7, within the listed assumptions.", "Primary-source verification, full source compatibility, dynamical stability and control, and operational performance remain open; cycle 1 does not establish them."]

\subsection{Research excerpt 5}
\textit{Source and field: }artifact:745ebe9c-6fd0-4f37-8dc9-007be29c954c:027a5b69-a96f-43f7-91ce-60af8ce26b57: skeptic\_output/reasoning\_assessments\par
[\{"step\_number": 1, "disposition": "accepted", "rationale": "Accepted as an explicit mathematical premise, not a realizability result. Differentiating the stated polynomial gives the stated velocity, whose maximum occurs at s=1/2 and equals 35D/(16T). The compact profile is smooth at its support boundary. The joined X is $C^{3}$, making v and the metric at least $C^{2}$.", "objection": "", "basis\_artifact\_ids": []\}, \{"step\_number": 2, "disposition": "accepted", "rationale": "Since X stays between 0 and D and L\textgreater{}R, neither laboratory intersects the support. The metric is Minkowski on their worldlines, and globally Minkowski before motion begins. Thus the initial synchronization and later proper-clock readings are consistent.

Euclidean slice distance is correctly distinguished from later radar distance.", "objection": "", "basis\_artifact\_ids": []\}, \{"step\_number": 3, "disposition": "accepted", "rationale": "The two flat legs each last L/u, and the center segment lasts T. Their endpoints match the stated center events. The curve is piecewise timelike, and the text explicitly excludes finite-acceleration implementation at its velocity corners.", "objection": "", "basis\_artifact\_ids": []\}, \{"step\_number": 4, "disposition": "accepted", "rationale": "The orthonormal coframe contains c dt and dx-vf dt, yielding the measured velocity dx/dt-vf. Substitution of the center path gives zero local velocity and dτ=dt.

Adding the two flat-leg proper times yields the stated total. The null-speed statement follows from the same coframe.", "objection": "", "basis\_artifact\_ids": []\}, \{"step\_number": 5,…

\subsection{Research excerpt 6}
\textit{Source and field: }session:745ebe9c-6fd0-4f37-8dc9-007be29c954c: strategy\_runs/0/pending\_investigation/reasoning\_steps\par
[\{"title": "Specify the geometry and the observer-dependent test", "argument": "Use coordinates x\textasciicircum{}0=ct and metric $ds^{2}$=-(dx\textasciicircum{}0)²+[dx-b f(r)dx\textasciicircum{}0]²+$dy^{2}$+$dz^{2}$, where $r^{2}$=(x-b x\textasciicircum{}0)²+$y^{2}$+$z^{2}$ and b is constant. Take R\textgreater{}0, $\Delta$\textgreater{}0, f=1 for r$\leq$R, f=0 for r$\geq$R+$\Delta$, with a smooth transition. The lapse is one, the spatial metric is δ\_ij, and the shift is $\beta$\textasciicircum{}i=(-bf,0,0). The future unit normal to a spatial slice is n=∂\_0+bf∂\_x. Define ε=T\_ab n\textasciicircum{}a n\textasciicircum{}b in energy-density units.

The weak energy condition requires T\_ab u\textasciicircum{}a u\textasciicircum{}b$\geq$0 for every timelike u, so ε\textless{}0 for this one observer is already decisive.", "epistemic\_status": "deduction", "basis\_artifact\_ids": []\}, \{"title": "Derive the necessary source from the Hamiltonian constraint", "argument": "With either overall sign convention for extrinsic curvature, the quadratic combination below is unchanged. Choose K\_ij=(∂\_i$\beta$\_j+∂\_j$\beta$\_i)/2. Its nonzero entries are K\_xx=-b f\_x, K\_xy=K\_yx=-b f\_y/2, and K\_xz=K\_zx=-b f\_z/2. Thus $K^{2}$=$b^{2}$$f_x^{2}$ and K\_ijK\textasciicircum{}ij=$b^{2}$[$f_x^{2}$+($f_y^{2}$+$f_z^{2}$)/2]. The spatial scalar curvature is zero. The Einstein constraint 16$\pi$G ε/c\textasciicircum{}4 = (3)R+$K^{2}$-K\_ijK\textasciicircum{}ij therefore gives ε=-c\textasciicircum{}4 $b^{2}$($f_y^{2}$+$f_z^{2}$)/(32$\pi$G).

This is a necessary component of the full source tensor, not a proposed matter model.", "epistemic\_status": "deduction", "basis\_artifact\_ids": []\}, \{"title": "Establish the obstruction at subluminal speed", "argument": "For a radial profile, $f_y^{2}$+$f_z^{2}$=f′(r)² $sin^{2}$$\theta$, with $\theta$ measured from the direction of translation. A smooth function changing from one to zero has f′$\neq$0 somewhere in the wall, and hence ε\textless{}0 on an open set away from the translation axis whenever b$\neq$0. No…

11 further research records are retained in the archived source snapshot.
\section*{Evidence and authorship note}
Snapshot SHA-256: d7930d1ec66a8c856c37f1fe483d06625ff95b2631508864d5444337ddc476a4
Generated by an AI author from the cited saved research snapshot. This document is not independent verification. No external bibliography was added.
\section*{Validation warnings}
\begin{itemize}
\item Evidence indexing omitted 1152 material saved research field fragment(s). Those fragments were not reviewed for this paper; the interpretation is provisional where the omitted evidence may matter. Consult the frozen source snapshot and omission ledger for full details.
\item Claim identifiers in abstract\_claim\_ids, problem\_framing\_claim\_ids, and methods\_claim\_ids require normalization: these arrays currently use zero-padded identifiers, whereas the claim register uses unpadded identifiers.
\item Several section-purpose strings are incomplete and require completion before this plan can guide a full draft; in particular, the source-bound proof, null-arrival comparison, requirements discussion, and quantum-source testing outline must retain their logical bridges.
\item Analytic-result claims are labeled conjecture to preserve their status as unverified conditional derivations, not to suggest that numerical experiments were performed or that independent proof certification exists.
\item The numerical illustrations in the evidence use different speeds: \$\textbackslash{}beta=0.1\$ gives a scale near \$1.1\textbackslash{}times10\textasciicircum{}\{44\}\textbackslash{},\textbackslash{}mathrm\{J\}\$, whereas \$\textbackslash{}beta=0.01\$ gives a scale near \$1.11\textbackslash{}times10\textasciicircum{}\{42\}\textbackslash{},\textbackslash{}mathrm\{J\}\$ for the same radii. Their factor-of-one-hundred difference follows the stated quadratic speed dependence and is not a contradiction.
\item The weak-energy-condition exclusion must not be upgraded to a demonstrated violation of every energy condition, a prohibition on quantum sourcing, or a theorem covering all subluminal warp-like metrics.
\item The temporal-function result applies to the stated globally defined metric; it does not establish causal safety of arbitrary formation protocols, multi-bubble configurations, or altered spacetime identifications.
\item The selected axial photon has not been established as the globally earliest signal. The traveler comparison does not require that stronger assertion.
\item No complete tensor, realizable source model, formation calculation, stability analysis, finite-acceleration transfer, or resource-matched mission comparison is supplied. Proposed follow-up calculations must remain clearly distinguished from completed results.
\item The evidence index explicitly omitted 1152 saved field fragment(s); see the generation coverage ledger and review the complete archived snapshot.
\item Structured paper indexing selected 434 source-linked text fragments from 7 archived sources without AI extraction calls. The complete frozen snapshot remains archived; this bounded index does not constitute review of every saved field.
\item C2 is labeled established, but cited saved evidence includes unverified, incomplete, or conjectural material.
\item C2 is labeled established while the frozen evidence contains counterexample or contradiction language; inspect the cited sources before accepting it.
\item C5 is labeled established, but cited saved evidence includes unverified, incomplete, or conjectural material.
\item C5 is labeled established while the frozen evidence contains counterexample or contradiction language; inspect the cited sources before accepting it.
\item C6 is labeled established, but cited saved evidence includes unverified, incomplete, or conjectural material.
\item C6 is labeled established while the frozen evidence contains counterexample or contradiction language; inspect the cited sources before accepting it.
\item The frozen evidence contains open proof/research obligation(s) requiring reconciliation: 60 item(s) are not listed by the author and are not presumed resolved; 10 author-listed item(s) have no exact frozen-record match. See the obligation reconciliation.
\item AI-authored analysis is not independent formal verification; model agreement does not establish proof.
\end{itemize}
\end{document}
